\documentclass[12pt]{article}

% ---- core math ----
% amsmath/amssymb must precede cleveref (loaded below); the canonical macro block
% re-declares them verbatim later, which is an idempotent no-op.
\usepackage{amsmath,amssymb,amsthm}

% ---- diagrams ----
\usepackage{tikz}
\usepackage{tikz-cd}

% ---- references / links ----
\usepackage{hyperref}
\usepackage{cleveref}

% ---- graphics / layout ----
\usepackage{graphicx}
\usepackage[margin=1in]{geometry}

% ---- tables ----
\usepackage{array}
\usepackage{booktabs}
\newcolumntype{L}[1]{>{\raggedright\arraybackslash}p{#1}}  % ragged-right fixed-width column

% ---- GrokRxiv sidebar ----
\usepackage{xcolor}

% ==== topological-phases-of-matter : canonical macros (v1, 2026-07-14) ====
% Requires amsmath, amssymb, amsfonts (mathfrak, mathbf, mathcal all standard).
\usepackage{amsmath,amssymb,amsfonts}

% --- condensation / probes ---
\newcommand{\cond}[1]{\underline{#1}}          % condensation X |-> \cond{X}
\newcommand{\cg}[1]{\underline{#1}}            % condensed group (alias, for readability)
\newcommand{\Cont}{\operatorname{Cont}}        % Cont(S,X)
\newcommand{\Shape}{\operatorname{Shape}}      % shape / homotopy type

% --- interaction space & locality ---
\newcommand{\Int}{\mathcal{I}}                 % \Int_{d,G} interaction space
\newcommand{\Ffun}{F}                          % F-function (interaction decay)
\newcommand{\BF}{\mathcal{B}_{F}}              % interaction Banach space for F

% --- moduli stacks ---
\newcommand{\Ham}{\mathfrak{Ham}}              % \Ham_{d,G}
\newcommand{\Gap}{\mathfrak{Gap}}              % \Gap_{d,G}
\newcommand{\Phase}{\mathfrak{Phase}}          % \Phase_{d,G}
\newcommand{\Phases}{\operatorname{Phases}}    % \Phases_{d,G} = pi_0 Shape Phase
\newcommand{\Weq}{\mathcal{W}}                 % class of equivalences to invert
\newcommand{\st}{\mathrm{st}}                  % stabilization superscript
\newcommand{\hfix}[1]{h#1}                     % homotopy-fixed-point superscript, e.g. \Phase_d^{\hfix{\cg{G}}}
\newcommand{\quotstack}[2]{[#1/#2]}            % action stack [Y/\cg{G}]

% --- gap ---
\newcommand{\gapof}[1]{\operatorname{gap}(#1)} % gap(H)
\newcommand{\gapbound}{\Delta}                  % uniform gap bound

% --- stacking / spectrum ---
% \boxtimes is built in (stacking product)
\newcommand{\IPcond}{\mathbf{IP}^{\mathrm{cond}}}  % invertible condensed phase spectrum

% --- transitions / discriminant / charge ---
\newcommand{\Sig}{\Sigma}                       % \Sig_f gapless discriminant
\newcommand{\rc}{\partial\nu}                   % relative transition charge in E^{q+1}(B,U_f)

% --- observables & solid K-theory ---
\newcommand{\Kop}{K_{\mathrm{op}}}             % operator / topological K-theory
\newcommand{\Kalg}{K_{\mathrm{alg}}}           % algebraic K-theory
\newcommand{\Solid}{\operatorname{Solid}}      % solidification functor

% --- disorder ---
\newcommand{\dis}{\Omega}                       % disorder hull \Omega = Q^{Z^d} (Q = local-configuration alphabet; NOT the F-function)
% ==== end canonical macros ====

% ==== Part V local notation (appended to coordination/notation.md symbol table) ====
\newcommand{\real}{\rho}                         % comparison / realizability map on pi_0
\newcommand{\SRE}{\mathrm{SRE}}                  % short-range-entangled sub-spectrum
\newcommand{\CP}{\mathrm{CP}}                    % commuting-projector qualifier
\newcommand{\Maj}{\mathcal{M}}                   % Majorana number (Kitaev chain Z2 invariant)
\newcommand{\winding}{w}                         % winding number
\newcommand{\Ztwo}{\mathbb{Z}_2}
\newcommand{\Zint}{\mathbb{Z}}
\newcommand{\Circ}{\mathbb{T}}                   % circle group U(1)=T (Ogata index target)
\newcommand{\im}{\operatorname{im}}
\DeclareMathOperator{\Pf}{Pf}
\DeclareMathOperator{\sgn}{sgn}
% ==== end local notation ====

% ---- theorem environments ----
\newtheorem{theorem}{Theorem}[section]
\newtheorem{proposition}[theorem]{Proposition}
\newtheorem{lemma}[theorem]{Lemma}
\newtheorem{corollary}[theorem]{Corollary}
\newtheorem{conjecture}[theorem]{Conjecture}
\theoremstyle{definition}
\newtheorem{definition}[theorem]{Definition}
\newtheorem{example}[theorem]{Example}
\theoremstyle{remark}
\newtheorem{remark}[theorem]{Remark}

% ---- GrokRxiv DOI sidebar (page 1 only) ----
\definecolor{grokgray}{RGB}{110,110,110}
\AddToHook{shipout/background}{%
  \ifnum\value{page}=1
    \begin{tikzpicture}[remember picture, overlay]
      \node[
        rotate=90,
        anchor=south,
        font=\Large\sffamily\bfseries\color{grokgray},
        inner sep=0pt
      ] at ([xshift=38pt, yshift=0.52\paperheight]current page.south west)
      {GrokRxiv:2026.07.bordism-realizability\quad
       [\,cond-mat.str-el\,]\quad
       14 Jul 2026};
    \end{tikzpicture}
  \fi
}

\hypersetup{
  colorlinks=true,
  linkcolor=blue!55!black,
  citecolor=green!45!black,
  urlcolor=blue!60!black
}

% Absorb the occasional stubborn line rather than overflow the margin.
\emergencystretch=2em

\title{\bfseries Physical Realizability of Bordism and Homotopy Classes\\ by Gapped Lattice Systems}

\author{Matthew Long \\
\textit{The YonedaAI Collaboration} \\
\textit{YonedaAI Research Collective} \\
Chicago, IL \\
\texttt{matthew@yonedaai.com} $\cdot$ \url{https://yonedaai.com}}

\date{14 July 2026\\[4pt]\small Part V of a six-part series}

\begin{document}
\maketitle

\begin{abstract}
\noindent
The invertible condensed phase spectrum $\IPcond_{d,G}$ of Part~IV assigns to each
short-range-entangled phase an abstract class in a bordism- or homotopy-theoretic
classification. The present paper asks the converse question: which of those classes are
actually produced by an explicit uniformly gapped lattice Hamiltonian? We formalize
realizability as the essential surjectivity, on $\pi_0$, of the comparison map
$\real$ from the stabilized stack of lattice models to the abstract spectrum, and we
separate two nested notions: realizability by \emph{some} gapped local Hamiltonian, and
realizability by a local \emph{commuting-projector} model. On the positive side we record,
in the condensed language, three facts that are theorems today: every group-cohomology
class in $H^{d+1}(G,U(1))$ is realized by an explicit $G$-symmetric commuting-projector
model carrying the prescribed edge anomaly (Chen--Gu--Liu--Wen; Else--Nayak); in spatial
dimension one Ogata's operator-algebraic index is a \emph{complete} invariant of
symmetry-protected phases with on-site finite symmetry, so the realized set coincides with the
invariant set there, while in dimension two the same index is known to exist and every
in-cohomology class is realized but its completeness is open; and commuting-projector
realizations face a hard obstruction: no local commuting-projector Hamiltonian carries a
nonzero electric Hall conductance (Kapustin--Fidkowski, under $U(1)$) or a nonzero chiral
central charge (Kapustin--Spodyneiko), so every chiral invertible class lies outside the
commuting-projector image. We assemble the known
constructions (group-cohomology models, Walker--Wang models for beyond-cohomology classes
in $3+1$ dimensions, the $E_8$ state, the Kitaev Majorana chain, cluster states) and the
known obstructions into a dimension-by-symmetry status table separating the torsion part
from the free part. We then state three numbered conjectures on the realizability image:
that it is exactly the short-range-entangled subspectrum; that chiral classes are
realizable by non-commuting gapped Hamiltonians though never by commuting projectors; and
that realizability descends over profinite disorder hulls $\dis=Q^{\Zint^d}$ precisely when
a condensed-cohomological obstruction vanishes. Accompanying Haskell computes the
Kitaev-chain $\Ztwo$ invariant across the parameter plane and the cluster-state string
order, verifying quantization, phase-boundary detection, and degradation under a
symmetry-breaking perturbation.
\end{abstract}

\tableofcontents

\section{Introduction}
\label{sec:intro}

\subsection{The realizability gap}

A classification theorem for topological phases produces an abstract invariant: a bordism
group, a stable-homotopy group, a value of some generalized cohomology theory. Part~IV of
this series constructs the target of such a classification inside the condensed-mathematics
program: the invertible condensed phase spectrum $\IPcond_{d,G}$, a connective condensed
spectrum whose homotopy groups are meant to enumerate the invertible gapped phases in
spatial dimension $d$ with symmetry $G$. What a spectrum enumerates and what a laboratory,
or a finite-dimensional local Hilbert space, can produce are not the same thing. A class
$c\in\pi_0\IPcond_{d,G}$ is an equivalence class of abstract data; to say that $c$ is
\emph{physically realizable} is to exhibit a concrete uniformly gapped local Hamiltonian
whose phase maps to $c$. The prospectus for this program names the discrepancy plainly:
the gap between abstract homotopy-TQFT classifications and explicit local lattice
realizations is a recognized issue, and not every abstract invertible field theory is known
to arise from a microscopic gapped system.

This paper is about that gap. We take the classification of Part~IV as given and study the
image of the comparison map
\[
  \real \colon \pi_0\Shape\bigl(\Phase_{d,G}\bigr) \longrightarrow \pi_0\IPcond_{d,G},
\]
from the set of phases of genuine lattice systems (components of the stabilized condensed
stack of Part~III) to the abstract spectrum. Realizability is the question of which
classes lie in $\im\real$; more precisely, since a phase is by definition already a class of
lattice systems, realizability is the essential surjectivity of $\real$ onto a
physically-meaningful subspectrum. We will argue that the honest present-day answer has
three parts: a large \emph{constructive} region where explicit models are known; a sharp
\emph{obstructed} region governed by a no-go theorem; and a residual \emph{open} region.

\subsection{Two notions of realizability}

The single most important distinction in this paper is between two nested classes of
Hamiltonians. Say that a class $c$ is
\begin{itemize}
  \item \emph{gapped-realizable} if some uniformly gapped local (or quasi-local) Hamiltonian
        has phase $c$; and
  \item \emph{commuting-projector realizable} (\emph{$\CP$-realizable}) if some Hamiltonian
        $H=\sum_x P_x$ with each $P_x$ a finite-range projector, $[P_x,P_y]=0$, has phase
        $c$.
\end{itemize}
Commuting-projector models are the exactly-solvable backbone of the field: they have exact
ground states, exact gaps, string operators, and a transparent entanglement structure. Every
$\CP$-realizable class is gapped-realizable, but (and this is the crux) the converse
fails, and it fails for a reason that is a rigorous theorem rather than a modeling
inconvenience. The Kapustin--Fidkowski theorem \cite{kapustin-fidkowski} states that a local
commuting-projector Hamiltonian has vanishing Hall conductance. A Chern insulator, an
integer quantum Hall state, the $E_8$ state (all manifestly gapped-realizable, all with
nonzero Hall response) are therefore \emph{not} $\CP$-realizable. Conflating the two
notions is the characteristic error this subject invites, and much of what follows is an
effort to keep them apart.

\subsection{What is proved, what is conjectured}

In keeping with the stance of the whole series, we label a statement a Theorem or
Proposition only when it follows from cited present-day results, and everything else is a
numbered Conjecture. The rigorous content of this paper is:
\begin{itemize}
  \item \textbf{Theorem~V-A} (\Cref{thm:VA}): every class in $H^{d+1}(G,U(1))$ is
        $\CP$-realizable, by the Chen--Gu--Liu--Wen fixed-point construction
        \cite{cglw-cohomology} with the Else--Nayak edge characterization \cite{else-nayak}.
  \item \textbf{Theorem~V-B} (\Cref{thm:VB}): in $d=1$, for on-site finite symmetry, Ogata's
        operator-algebraic index \cite{ogata-spt-chains,ogata-classification-review} is a
        complete invariant, so realized $=$ invariant there; in $d=2$ the $H^3$ index exists
        \cite{ogata-h3-index} and every in-cohomology class is realized, but completeness is
        open (\Cref{rem:VB2d}).
  \item \textbf{Theorem~V-C} (\Cref{thm:VC}): the commuting-projector no-go, zero electric
        Hall conductance under $U(1)$ (Kapustin--Fidkowski \cite{kapustin-fidkowski}) and zero
        chiral central charge (Kapustin--Spodyneiko \cite{kapustin-spodyneiko-thermal}),
        exactly bounds the $\CP$-image.
  \item Two elementary Propositions with complete proofs: the Kitaev-chain winding invariant
        (\Cref{prop:kitaev}) and the cluster-state string order (\Cref{prop:cluster}), both
        of which are also checked by the accompanying Haskell.
\end{itemize}
The conjectural content is three numbered statements about the global shape of $\im\real$
(\Cref{sec:image}): that it is the short-range-entangled subspectrum (Conjecture~V-1); that
chiral classes are gapped-realizable but never $\CP$-realizable (Conjecture~V-2); and that
realizability descends over profinite disorder hulls under a condensed-cohomological
obstruction (Conjecture~V-3).

\subsection{Relation to companion papers}
\label{sec:companions}

This is Part~V of a six-part program that recasts topological phases of matter in the
condensed-mathematics paradigm of Clausen--Scholze \cite{clausen-scholze-condensed}. The
parts compose as a modular hierarchy, each taking the previous as input.

Part~I, \emph{Condensed Locality} \cite{paperLocality}, makes the interaction space a Banach
space $\BF$ and promotes the Heisenberg dynamics to a morphism of condensed objects; it
supplies the Lieb--Robinson control that gives meaning to a ``uniformly gapped family'' over
a profinite base. Part~II, \emph{Positivity and Condensed State Spaces}
\cite{paperPositivity}, attaches the quasi-local $C^*$-algebra and, through Aoki's
solidification bridge \cite{aoki-solidification}, a solid-$K$-theory invariant; the
operator-$K$ invariants that
appear in our status table (Hall conductance, Chern number) are values of that functor.
Part~III, \emph{The Uniformly Gapped Substack} \cite{paperGap}, cuts out the substack
$\Gap_{d,G}$ whose stabilized $\pi_0$ is the set $\Phases_{d,G}$ that our map $\real$ starts
from; the robustness of that $\pi_0$ under perturbation is what makes ``the phase of a
lattice model'' well defined. Part~IV, \emph{From Lattice Models to Effective Field
Theories} \cite{paperEFT}, group-completes the invertible sector into the target spectrum
$\IPcond_{d,G}$; realizability is precisely the surjectivity question for the comparison map
that Part~IV constructs. Part~VI, the synthesis \cite{paperSynthesis}, assembles the five
modules into a single program and records the Master Conjecture; our obstruction subgroup is
one of the two hard walls (the other being the undecidability of the gap in Part~III) that
bound the synthesis.

The reader needs from Parts~III--IV only the following interface, restated in
\Cref{sec:framework}: a phase is a component of a stabilized condensed stack, stacking gives
it an abelian-group structure, and $\IPcond_{d,G}$ is the invertible part group-completed.
No result of this paper depends on the conjectural global stack structure; we use the
module-level theorems of the companions and cite them where needed.

\subsection{Outline}

\Cref{sec:framework} sets up the comparison map and the two realizability notions inside the
condensed program. \Cref{sec:cohomology,sec:beyond} survey positive constructions: the
group-cohomology models and the cluster state, then beyond-cohomology and topological-order
models together with the Kitaev Majorana chain. \Cref{sec:complete} states the low-dimension
completeness theorem and is careful about its scope. \Cref{sec:nogo} states the
Kapustin--Fidkowski obstruction and draws the $\CP$-versus-gapped line. \Cref{sec:table}
gathers everything into a status table. \Cref{sec:image} states the conjectures on the
realizability image, and \Cref{sec:disorder} treats disorder-robust realizability over
profinite hulls. \Cref{sec:haskell} describes the computational checks.
\Cref{sec:discussion,sec:conclusion} discuss limitations and conclude.

\section{The realizability question in the condensed program}
\label{sec:framework}

\subsection{The interface from Parts III and IV}

We recall only what we need. Fix a spatial dimension $d$, a lattice $L$, on-site Hilbert
spaces, an $F$-function locality class, and a symmetry group $G$ with condensation $\cg{G}$.
Part~I organizes the admissible $G$-symmetric interactions into the condensed moduli object
$\Ham_{d,G}$, whose value on a profinite probe $S$ is the set of $S$-continuous families of
interactions. Part~III singles out the \emph{uniformly} gapped substack
\[
  \Gap_{d,G}(S) \;=\; \bigcup_{\gapbound>0}\Bigl\{\,H\in\Ham_{d,G}(S) : \inf_{s\in S}\gapof{H_s}\ge\gapbound \,\Bigr\},
\]
where the word \emph{uniformly} (the infimum over the probe) is essential: a family
with $\inf_s\gapof{H_s}=0$ is not a member. Inverting the class $\Weq$ of gapped adiabatic /
finite-depth equivalences \cite{bmns-automorphic} and stabilizing by product-state ancillas gives the phase
$\infty$-groupoid $\Phase_{d,G}=(\Gap_{d,G}[\Weq^{-1}])^{\st}$, and the set of phases is
\[
  \Phases_{d,G} \;=\; \pi_0\Shape\bigl(\Phase_{d,G}\bigr).
\]
Stacking $\boxtimes$ makes $\Phase_{d,G}$ symmetric monoidal; its invertible objects form a
Picard $\infty$-groupoid, and Part~IV group-completes them into the connective condensed
spectrum $\IPcond_{d,G}$. Writing $\Phases^{\times}_{d,G}\subseteq\Phases_{d,G}$ for the
invertible (short-range-entangled) phases, group completion is a map
\[
  \gamma \colon \Phases^{\times}_{d,G} \longrightarrow \pi_0\IPcond_{d,G}.
\]

\subsection{The comparison map and realizability}

Part~IV's central comparison is between the microscopic invariant computed on the lattice
and the deformation class of the associated effective field theory. Composing the lattice
phase with that comparison gives a map to the abstract bordism/homotopy classification;
denote its value on $\pi_0$ by
\[
  \real \colon \Phases_{d,G} \longrightarrow \mathcal{C}_{d,G},
\]
where $\mathcal{C}_{d,G}$ is the abstract classification group of the relevant EFT: for
invertible phases, $\pi_0\IPcond_{d,G}$ after realification, conjecturally a connective cover
of Kubota's $\Omega$-spectrum \cite{kubota-omega-spectrum} or the Freed--Hopkins bordism
target \cite{freed-hopkins}. The direction of $\real$ matters. Constructing a lattice model
and reading off its abstract class is the ``forward'' direction and is always available.
Realizability is the ``backward'' question: given $c\in\mathcal{C}_{d,G}$, is there a phase
$P\in\Phases_{d,G}$ with $\real(P)=c$?

\begin{definition}[Realizability]
\label{def:realizable}
A class $c\in\mathcal{C}_{d,G}$ is \emph{gapped-realizable} if $c\in\im\real$, i.e.\ there is
a uniformly gapped $G$-symmetric local Hamiltonian whose phase $P$ satisfies $\real(P)=c$. It
is \emph{$\CP$-realizable} if such an $H$ can moreover be taken to be a local
commuting-projector Hamiltonian $H=\sum_x P_x$, $P_x^2=P_x=P_x^\ast$, $[P_x,P_y]=0$, each
$P_x$ of finite range. Write
\[
  \im\real \;\supseteq\; \im\real^{\CP}
\]
for the two images; $\im\real^{\CP}\subseteq\im\real$ always.
\end{definition}

The whole paper is the study of the two subsets $\im\real^{\CP}\subseteq\im\real\subseteq
\mathcal{C}_{d,G}$. We record the reformulation that ties realizability to the language of
the program.

\begin{proposition}[Realizability as essential surjectivity]
\label{prop:esssurj}
$c$ is gapped-realizable if and only if the comparison functor
$\Phase_{d,G}\to\mathcal{C}_{d,G}$ (viewing $\mathcal{C}_{d,G}$ as a discrete groupoid) hits
the object $c$ up to isomorphism; that is, $c\in\im\real$ iff $c$ is in the essential image
on $\pi_0$. Equivalently, gapped-realizability of all of $\mathcal{C}_{d,G}$ is the
statement that $\real$ is surjective.
\end{proposition}

\begin{proof}
Immediate from the definitions: $\pi_0$ of the comparison functor is the map $\real$ of
sets, and a set map hits $c$ iff $c$ is in its image, which for a functor into a discrete
groupoid is the essential image on objects. There is no content beyond unwinding the
notation; the content is entirely in \emph{which} $c$ are hit, addressed in the following
sections.
\end{proof}

\begin{remark}[Why $\CP$ is not a technicality]
\label{rem:cp}
One might hope that $\CP$-realizability and gapped-realizability coincide up to stable
equivalence, so that the exactly-solvable models exhaust the phases. They do not:
\Cref{thm:VC} exhibits a permanent obstruction. The correct picture is that
$\im\real^{\CP}$ is a proper subgroup of $\im\real$ whenever chiral phases exist (i.e.\ for
$d=2$ without antiunitary protection, and in higher dimensions with the appropriate
gravitational response). Keeping the two apart is what lets us state positive results
(constructions are usually $\CP$) and the obstruction (which is about $\CP$) without
contradiction.
\end{remark}

\subsection{Short-range entanglement}

The invertible phases are the short-range-entangled ones \cite{freed-sre}: a phase $P$ is invertible for
$\boxtimes$ iff there is $P'$ with $P\boxtimes P'$ trivial (a product state) after
stabilization. We write $\SRE_{d,G}\subseteq\Phases_{d,G}$ for this subset; group completion
$\gamma$ restricts to it. Non-invertible topological order (nonzero total quantum dimension,
anyons that cannot be condensed away) lies outside $\SRE$ and outside the spectrum-level
classification altogether; it needs the higher-categorical machinery flagged in Part~VI and
is not classified by $\IPcond_{d,G}$. Our realizability question for the \emph{spectrum} is
therefore a question about $\SRE$; topologically ordered models such as the toric code enter
only as $\CP$-realizations of \emph{trivial} invertible class (they are invertible-trivial
but intrinsically ordered), and the chiral topological orders (e.g.\ Kitaev's non-abelian
phase) enter through their invertible \emph{edge} data, the chiral central charge.

\section[Positive realizability I: group-cohomology SPT phases]{Positive realizability I:\\ group-cohomology SPT phases}
\label{sec:cohomology}

\subsection{The Chen--Gu--Liu--Wen construction}

The first and largest constructive region is the group-cohomology part of the
symmetry-protected topological (SPT) classification. Let $G$ be a finite group acting
on-site and unitarily; work in spatial dimension $d$, so the relevant cocycle degree is
$d+1$.

\begin{theorem}[Cohomological classes are realized; V-A]
\label{thm:VA}
For every finite group $G$ and every class $\omega\in H^{d+1}(G,U(1))$ there is an explicit
$G$-symmetric local commuting-projector Hamiltonian $H_\omega$ on a lattice with on-site
$G$-representation, such that:
\begin{enumerate}
  \item[(i)] $H_\omega$ has a unique gapped ground state on any closed $d$-manifold; the
        ground state is a renormalization-group fixed-point wavefunction and is
        short-range-entangled (preparable from a product state by a finite-depth circuit once
        the on-site $G$-symmetry constraint is dropped --- no $G$-symmetric finite-depth
        circuit prepares it, which is exactly what makes the phase a nontrivial SPT);
  \item[(ii)] the phase of $H_\omega$ is trivial as an ungauged state but nontrivial as a
        $G$-SPT, and the assignment $\omega\mapsto[H_\omega]$ is a group homomorphism
        $H^{d+1}(G,U(1))\to\SRE_{d,G}$ that is injective on the group-cohomology subgroup;
  \item[(iii)] the boundary of $H_\omega$ carries an anomalous $G$-action whose obstruction
        to being realized on-site is exactly $\omega$.
\end{enumerate}
In particular every group-cohomology class is $\CP$-realizable.
\end{theorem}

\begin{proof}
This is the content of \cite{cglw-cohomology}, with the edge statement (iii) supplied by
\cite{else-nayak}; we give the structure of the argument and cite the theorems for the
analytic steps. Chen--Gu--Liu--Wen construct, from a representative inhomogeneous
$(d+1)$-cocycle $\nu_{d+1}\colon G^{d+1}\to U(1)$, a fixed-point wavefunction on a
triangulated $d$-manifold: to each branched triangulation one assigns an amplitude that is a
product of cocycle phases $\nu_{d+1}$ over the $(d+1)$-simplices of a filling, and the
cocycle condition guarantees that the amplitude is independent of the filling and defines a
short-range-correlated state. The parent Hamiltonian is a sum of projectors, one per local
patch, each projecting onto the local fixed-point subspace; because the fixed-point
conditions on overlapping patches are compatible, the projectors commute, giving a
commuting-projector Hamiltonian with the wavefunction as its unique frustration-free ground
state. That establishes (i). The map on classes is additive because stacking two decorated
wavefunctions multiplies their cocycle phases, i.e.\ adds the cohomology classes; injectivity
on the group-cohomology subgroup is the statement that distinct classes give
non-adiabatically-connected models, proved by the edge obstruction of (iii). For (iii),
Else--Nayak show that the boundary of such an SPT carries a $G$-action that cannot be
extended to an on-site action of $G$ on the boundary Hilbert space alone; the obstruction to
on-site-ness lives in $H^{d+1}(G,U(1))$ and equals $\omega$. This anomalous edge action is
what protects the bulk phase and certifies nontriviality.
\end{proof}

\begin{remark}[Scope of Theorem~V-A]
\label{rem:VAscope}
Theorem~V-A covers exactly the \emph{in-cohomology} subgroup of the SPT classification. It
does \emph{not} claim to realize classes outside $H^{d+1}(G,U(1))$: the beyond-group-cohomology
classes of Kapustin \cite{kapustin-cobordism}, detected by cobordism invariants rather than
group cohomology, require different constructions (\Cref{sec:beyond}). Nor does on-site
finiteness extend for free to continuous $G$ or to spatial symmetry; those need their own
constructions and are not asserted here.
\end{remark}

\subsection{The cluster state as a worked example}

The one-dimensional cluster state is the smallest nontrivial instance of \Cref{thm:VA} and
the model our Haskell computes with, so we give it explicitly. Take a chain of qubits with
the symmetry $G=\Ztwo\times\Ztwo$ generated by
$X_{\mathrm{even}}=\prod_{j\ \mathrm{even}}X_j$ and
$X_{\mathrm{odd}}=\prod_{j\ \mathrm{odd}}X_j$. The relevant cohomology is
$H^2(\Ztwo\times\Ztwo,U(1))\cong\Ztwo$, so there is exactly one nontrivial SPT class, and it
is realized by the cluster state.

\begin{definition}[Cluster state]
\label{def:cluster}
The cluster state $\lvert C\rangle$ on a chain of $n$ qubits is the unique (up to boundary
conditions) common $+1$ eigenstate of the commuting stabilizers
\[
  K_j \;=\; Z_{j-1}\,X_j\,Z_{j+1}, \qquad 1<j<n,
\]
together with boundary stabilizers. It is the frustration-free ground state of
$H_{\mathrm{cl}}=-\sum_j K_j$, a commuting-projector Hamiltonian (after $P_j=(1+K_j)/2$).
\end{definition}

The SPT order of $\lvert C\rangle$ is invisible to any local order parameter but visible to a
nonlocal string operator. The following is elementary and is verified exactly by the Haskell
stabilizer engine.

\begin{proposition}[Cluster-state string order]
\label{prop:cluster}
Fix $a<b$ with $b-a$ even, and set
\[
  S_{a,b} \;=\; Z_{a}\,\Bigl(\textstyle\prod_{r=0}^{(b-a)/2-1} X_{a+1+2r}\Bigr)\,Z_{b}
  \;=\; Z_a\,X_{a+1}X_{a+3}\cdots X_{b-1}\,Z_b .
\]
Then in the cluster state $\langle C\rvert S_{a,b}\lvert C\rangle = +1$ exactly, while every
single-site expectation vanishes, $\langle C\rvert Z_j\lvert C\rangle=\langle C\rvert
X_j\lvert C\rangle=0$. Moreover, under the symmetry-breaking on-site rotation
$U(\theta)=\prod_j e^{-i\theta Z_j/2}$ (which breaks $X_{\mathrm{even}},X_{\mathrm{odd}}$),
\[
  \langle C\rvert U(\theta)^\ast\,S_{a,b}\,U(\theta)\lvert C\rangle \;=\; (\cos\theta)^{m},
  \qquad m=\tfrac{b-a}{2} = \#\{\text{$X$ factors in }S_{a,b}\},
\]
so the string order degrades monotonically from $1$ at $\theta=0$ to $0$ at
$\theta=\pi/2$.
\end{proposition}

\begin{proof}
The stabilizer group $\mathcal{S}=\langle K_j\rangle$ is abelian and $\lvert C\rangle$ is its
unique $+1$ eigenvector, so for a Pauli operator $P$ one has $\langle C\rvert P\lvert
C\rangle=+1$ if $P\in\mathcal{S}$, $=-1$ if $-P\in\mathcal{S}$, and $=0$ otherwise (because
then $P$ anticommutes with some stabilizer $g$, whence $\langle C\rvert P\lvert
C\rangle=\langle C\rvert PggP^{-1}\cdots\rangle$ forces a sign flip and the value is its own
negative). The telescoping identity
\[
  K_{a+1}K_{a+3}\cdots K_{b-1}
  = \bigl(Z_a X_{a+1}Z_{a+2}\bigr)\bigl(Z_{a+2}X_{a+3}Z_{a+4}\bigr)\cdots
  = Z_a\,X_{a+1}X_{a+3}\cdots X_{b-1}\,Z_b ,
\]
where the internal $Z^2=1$ cancellations occur at the shared even-offset sites, shows that
this product of alternate stabilizers equals $S_{a,b}$ exactly. Hence $S_{a,b}\in\mathcal{S}$
and $\langle C\rvert S_{a,b}\lvert C\rangle=+1$. The single-site $Z_j$ anticommutes with $K_j$
(through the central $X_j$), and the single-site $X_j$ anticommutes with $K_{j\pm1}$ (through
their $Z_j$ factors); hence each has expectation $0$. No product of a single $Z$ or $X$ with a
stabilizer lands back in $\mathcal{S}$, so these are genuinely zero. For the rotated expectation, $U(\theta)^\ast Z_j U(\theta)=Z_j$ and
$U(\theta)^\ast X_j U(\theta)=\cos\theta\,X_j-\sin\theta\,Y_j$, so expanding
$U^\ast S_{a,b}U$ over the $m$ rotated $X$ factors gives $2^m$ Pauli terms. The all-$X$ term
is $(\cos\theta)^m S_{a,b}$ with expectation $(\cos\theta)^m$. Any other term differs from
$S_{a,b}$ by replacing at least one $X_k$ with $Y_k=iX_kZ_k$, i.e.\ equals a phase times
$S_{a,b}\cdot\prod_{k\in T}Z_k$ for a nonempty set $T$ of the $X$-sites; since the only
$X$-free element of $\mathcal{S}$ is the identity, $\prod_{k\in T}Z_k\notin\mathcal{S}$ for
$T\ne\varnothing$, so $S_{a,b}\prod_{k\in T}Z_k\notin\mathcal{S}$ (up to sign) and its
expectation is $0$. Thus all cross terms vanish and the total is $(\cos\theta)^m$. Monotone
decrease on $[0,\pi/2]$ is clear.
\end{proof}

\begin{remark}
\Cref{prop:cluster} is exactly the SPT-diagnostic behaviour: a nonlocal string order that is
quantized in the protected phase and destroyed only when the protecting symmetry is broken,
with no accompanying local order parameter. It realizes the nontrivial class of
$H^2(\Ztwo\times\Ztwo,U(1))\cong\Ztwo$, the smallest case of \Cref{thm:VA}. The Haskell of
\Cref{sec:haskell} computes $\langle S_{a,b}\rangle$ and its $\theta$-degradation from the
stabilizer group directly, with no closed-form shortcut.
\end{remark}

\section[Positive realizability II: beyond cohomology and topological order]{Positive realizability II:\\ beyond cohomology and topological order}
\label{sec:beyond}

\subsection{Beyond group cohomology: Walker--Wang models}

Not every invertible or topologically ordered phase is captured by group cohomology. In
$3+1$ dimensions there are SPT and symmetry-enriched phases whose invariants are cobordism
classes outside the image of group cohomology \cite{kapustin-cobordism}. A general
commuting-projector construction for a large class of these is available.

\begin{proposition}[Walker--Wang realizations]
\label{prop:ww}
Given a premodular (braided fusion) category $\mathcal{B}$, the Walker--Wang construction
\cite{walker-wang} produces a $3+1$-dimensional local commuting-projector Hamiltonian whose
ground state on a closed $3$-manifold is a fixed-point state, and whose $2$-dimensional
boundary carries the topological order associated with $\mathcal{B}$. When $\mathcal{B}$ is
modular the bulk is invertible (a $3+1$d SPT-like state) with a chiral boundary; when
$\mathcal{B}$ is a symmetric fusion category the construction reproduces group-cohomology and
certain beyond-cohomology SPTs.
\end{proposition}

\begin{proof}
Cited: \cite{walker-wang} construct the state sum and its commuting-projector parent
Hamiltonian from the premodular data, and identify the boundary theory. We use it only as a
source of $\CP$-realizations for classes beyond \Cref{thm:VA}; no new argument is supplied.
\end{proof}

\begin{remark}[The chiral boundary caveat]
\label{rem:wwchiral}
Walker--Wang models are $\CP$ in the $3+1$d bulk, but a modular input forces the $2+1$d
boundary to be chiral, and by \Cref{thm:VC} that boundary cannot itself be a $\CP$ model.
This is not a contradiction (the bulk is $\CP$, the boundary is not), but it is a clean
illustration of why $\CP$-realizability is a property of a specific model in a specific
dimension, not a phase-intrinsic and dimension-blind attribute.
\end{remark}

\subsection{Chiral invertible phases: the \texorpdfstring{$E_8$}{E8} state}

The cleanest bosonic example of the obstruction to come is the $E_8$ state: a $2+1$d bosonic
short-range-entangled phase with no symmetry, nontrivial because its edge carries eight
chiral bosons with chiral central charge $c_-=8$ and a thermal Hall response. It generates
the free part $\Zint$ of the $2+1$d bosonic invertible classification. The $E_8$ state is
gapped-realizable (it is a genuine gapped phase, described by a $K$-matrix Chern--Simons
theory and realizable in coupled-layer and network constructions), but because $c_-\ne0$ it
has a nonzero thermal Hall conductance, and \Cref{thm:VC}(ii) (Kapustin--Spodyneiko
\cite{kapustin-spodyneiko-thermal}) forbids any local commuting-projector realization. This
obstruction rests on the chiral central charge, \emph{not} on the electric Kapustin--Fidkowski
no-go: the bosonic $E_8$ state carries no $U(1)$ charge and has $\sigma_H=0$, so part (i) of
\Cref{thm:VC} does not apply to it and part (ii) is what does the work. It is the bosonic
analogue, at the level of gravitational response, of the fermionic integer quantum Hall state
\cite{tknn-1982}; its nonzero chiral
central charge is the edge datum that also governs chiral topological order more broadly
\cite{kitaev-honeycomb}.

\subsection{Fermionic \texorpdfstring{$\Ztwo$}{Z2}: the Kitaev Majorana chain}

The canonical fermionic example is the Kitaev chain \cite{kitaev-majorana-wire}, the $1+1$d
class-$D$/BDI superconductor whose two phases are distinguished by a $\Ztwo$ invariant and
whose topological phase binds an unpaired Majorana zero mode at each end. Its $\Ztwo$ label is
the class-$D$ entry of the periodic table of free-fermion phases \cite{kitaev-periodic-table},
whose systematic form is the operator-$K$-theoretic classification of topological insulators
and superconductors \cite{thiang-ktheory,prodan-sb} together with its bulk-edge correspondence
\cite{bourne-kellendonk-rennie,kubota-controlled}; the closely related Su--Schrieffer--Heeger
chain \cite{ssh-1979} is its chiral-symmetric cousin. It is the fermionic counterpart of the
cluster state, and our Haskell computes its invariant across the parameter plane.

Write the chain with spinless fermions $c_j$, hopping $t$, chemical potential $\mu$, and
$p$-wave pairing amplitude $\Delta_{\mathrm{p}}$ (written with a subscript to keep it visibly
distinct from the uniform gap bound $\gapbound=\Delta$ of \Cref{sec:framework}):
\[
  H = \sum_j\Bigl[-t\,(c_j^\dagger c_{j+1}+\mathrm{h.c.})
      -\mu\bigl(c_j^\dagger c_j-\tfrac12\bigr)
      +\Delta_{\mathrm{p}}\,(c_j c_{j+1}+\mathrm{h.c.})\Bigr].
\]
In momentum space the Bogoliubov--de Gennes Bloch Hamiltonian is
$h(k)=d_z(k)\,\tau_z+d_y(k)\,\tau_y$ in the Nambu basis, with
\[
  d_z(k) = -2t\cos k-\mu, \qquad d_y(k) = 2\Delta_{\mathrm{p}}\sin k, \qquad d_x(k)\equiv 0,
\]
where $\tau_{x,y,z}$ are Pauli matrices in particle--hole space and the $\tau_x$ component
$d_x$ vanishes identically because the pairing is real. The single-particle spectrum is
$\pm\sqrt{d_y^2+d_z^2}$, so the bulk gap closes exactly when $d_y=d_z=0$, i.e.\ at
$\sin k=0$ and $\mu=\mp2t$, i.e.\ on the discriminant $|\mu|=2|t|$ (for $\Delta_{\mathrm{p}}\ne0$).

\begin{proposition}[Kitaev-chain $\Ztwo$ invariant]
\label{prop:kitaev}
For $\Delta_{\mathrm{p}}\ne0$ and $|\mu|\ne2|t|$ the map $k\mapsto(d_y(k),d_z(k))\in\mathbb{R}^2\setminus\{0\}$,
$k\in[0,2\pi)$, has a well-defined winding number $\winding\in\Zint$ about the origin, and
\[
  \winding =
  \begin{cases}
    1, & |\mu|<2|t| \quad(\text{topological}),\\[2pt]
    0, & |\mu|>2|t| \quad(\text{trivial}).
  \end{cases}
\]
Equivalently the Majorana number $\Maj=\sgn\!\bigl(d_z(0)\bigr)\,\sgn\!\bigl(d_z(\pi)\bigr)
=\sgn\!\bigl((\mu+2t)(\mu-2t)\bigr)=\sgn(\mu^2-4t^2)$ equals $-1$ in the topological phase
and $+1$ in the trivial phase. The invariant is quantized on the gapped locus and jumps
precisely on $\Sig=\{|\mu|=2|t|\}$; the jump is the $1$d instance of the relative transition
charge $\rc$ of the program.
\end{proposition}

\begin{proof}
The curve $k\mapsto(d_y(k),d_z(k))=(2\Delta_{\mathrm{p}}\sin k,\,-2t\cos k-\mu)$ is an ellipse with
semi-axes $2|\Delta_{\mathrm{p}}|$ (horizontal) and $2|t|$ (vertical) centered at $(0,-\mu)$. For
$\Delta_{\mathrm{p}},t\ne0$ it is a nondegenerate closed curve; its winding number about the origin is
$1$ if the origin lies inside the ellipse and $0$ if outside (the ellipse is convex and
traversed once). The origin is inside iff $\mu^2/(2t)^2<1$, i.e.\ $|\mu|<2|t|$, giving the
stated $\winding$. The Majorana number is the product of the signs of $d_z$ at the two
particle--hole-invariant momenta $k=0,\pi$ (where $d_y=0$), a standard $\Ztwo$ Pfaffian
invariant for class $D$; here $d_z(0)=-2t-\mu$ and $d_z(\pi)=2t-\mu$, so
$\Maj=\sgn(-2t-\mu)\sgn(2t-\mu)=\sgn(\mu^2-4t^2)$. Both $\winding\bmod2$ and $\Maj$ change
value exactly across $|\mu|=2|t|$, where the curve passes through the origin and the bulk gap
closes; on the gapped locus each is locally constant and integer/sign quantized.
\end{proof}

\begin{remark}[What this realizes, and what it does not]
\label{rem:kitaevfree}
The Kitaev chain $\CP$-realizes nothing chiral: it is $1+1$d and its invariant is
torsion ($\Ztwo$), and at the fixed point $\mu=0,\ t=\Delta_{\mathrm{p}}$ it is itself a
commuting-(Majorana-)projector model. It illustrates the torsion, $\CP$-realizable corner of
the table. The free ($\Zint$) invariants live in even space dimension (Chern number,
$d=2$) and are governed by the obstruction of \Cref{sec:nogo}; the Kitaev chain does not
reach them.
\end{remark}

In two dimensions the fermionic story is richer and, importantly, also constructive: the
group-supercohomology models of Gu--Wen \cite{gu-wen-supercohomology} realize a large class of
interacting fermionic SPTs, the discrete-spin-structure construction of Tarantino--Fidkowski
\cite{tarantino-fidkowski} gives explicit commuting-projector models for \emph{all} $2+1$d
fermionic SPTs with symmetry $G_f=G\times\Ztwo^f$ (including all eight members of the
$\Zint_8$ classification for $G=\Ztwo$), and the general construction-and-classification
programme of Wang--Gu \cite{wang-gu-fermionic} extends this. These are torsion classes and are
$\CP$-realizable; they populate the fermionic rows of the status table (\Cref{tab:status}).

\subsection{Homotopy of the space of models}

Realizability is not only a $\pi_0$ question. The higher homotopy of the space of gapped
Hamiltonians records adiabatic pumps ($\pi_1$) and higher families ($\pi_n$); realizing a
class in $\pi_n$ means building an $n$-parameter family of gapped models with the prescribed
higher invariant; the $\pi_1$ generator is the adiabatic charge pump of Thouless
\cite{thouless-pump}, whose rigorous lattice avatar is the higher Berry class of
Kapustin--Sopenko \cite{kapustin-sopenko-noether}, and the homotopical foundations of such
parametrized families are set up by Beaudry et al.\ \cite{beaudry-etal}; a complementary
classification by the homotopy of the space of states of a $C^*$-algebra appears in
\cite{denittis-states,denittis-rendel-weyl}. For $2+1$d topological order,
Aasen--Wang--Hastings
\cite{aasen-wang-hastings} propose explicit generators of $\pi_1$, $\pi_2$, $\pi_3$ of the
space of Hamiltonians via automorphisms of topological order and honeycomb/automorphism
codes. We stress that these homotopy groups are \emph{conjectural}: \cite{aasen-wang-hastings}
give constructions and strong evidence but not a proof that these are the full homotopy
groups, and we cite them as such. The moduli-space perspective of Hsin--Wang
\cite{hsin-wang-moduli} is the closest ordinary-topology precedent for the object Part~IV
condenses.

\section{Completeness in low dimension}
\label{sec:complete}

In dimension one, realizability is completely settled for on-site finite symmetry, because the
operator-algebraic invariant is a \emph{complete} invariant: the realized set is forced to
equal the invariant set. In dimension two the same operator-algebraic index is known to exist
and every in-cohomology class is realized, but its completeness is open, so only the
surjectivity half survives there. The abstract target is the group-cohomology classification,
itself a special case of the generalized-cohomology description of SPT phases
\cite{gaiotto-jf,xiong-minimalist}.

\begin{theorem}[Completeness in one dimension; V-B]
\label{thm:VB}
Let $G$ be a finite group acting on-site on a quantum spin chain, and consider $G$-symmetric
gapped ground-state phases in the operator-algebraic (split/approximately-factorized) sense.
Then in $d=1$ the $H^2(G,U(1))$-valued index is a \emph{complete} invariant of the symmetric
phase \cite{ogata-spt-chains,ogata-classification-review}. Consequently, in $d=1$ with
on-site finite symmetry, the map $\real$ restricted to these SPT phases is a \emph{bijection}
onto $H^2(G,U(1))$: every class is realized (by \Cref{thm:VA}) and no two inequivalent models
share an index. Realizability is settled there: realized $=$ invariant.
\end{theorem}

\begin{proof}
The completeness statement is Ogata's theorem: the $d=1$ classification of $G$-symmetric
gapped phases by the second-cohomology index is proved in
\cite{ogata-spt-chains,ogata-classification-review}, for on-site finite $G$ in the stated
operator-algebraic framework. Surjectivity onto $H^2(G,U(1))$ is \Cref{thm:VA} (the CGLW
models realize every class). Injectivity of $\real$ on these phases is exactly the
completeness of the index. Composing, $\real$ is a bijection in $d=1$.
\end{proof}

\begin{remark}[Two dimensions: the index exists, completeness is open; V-B$'$]
\label{rem:VB2d}
In $d=2$ the situation is genuinely weaker, and it is important not to overstate it. Ogata
constructs an $H^3(G,\Circ)$-valued index for $G$-symmetric gapped phases with on-site finite
symmetry \cite{ogata-h3-index}, and proves it is a \emph{well-defined invariant}; combined
with the CGLW construction (\Cref{thm:VA}), the map $\real$ is therefore \emph{surjective}
onto the in-cohomology classes $H^3(G,\Circ)$, so every such class is realized. But whether
this index is a \emph{complete} invariant (whether it separates all $d=2$ SPT phases with
on-site finite symmetry) is \emph{not} established: Ogata's theorem gives existence and
invariance of the index, not completeness. We therefore record $d=2$ completeness as an
\textbf{open problem}, not a theorem, and do \emph{not} assert realized $=$ invariant in
$d=2$; only the surjectivity half (every in-cohomology class is realized) is settled there.
This matches the identical treatment in Part~IV~\cite{paperEFT}.
\end{remark}

\begin{remark}[Scope: do not extrapolate]
\label{rem:VBscope}
Theorem~V-B is a completeness statement only for \emph{on-site finite} symmetry in $d=1$; its
$d=2$ analogue is the open problem of \Cref{rem:VB2d} (the index exists but completeness is
unproven). It must not be read as a completeness statement for:
\begin{itemize}
  \item continuous symmetry groups (e.g.\ $U(1)$, $SU(2)$), where the index theory and the
        classification differ;
  \item spatial (crystalline) symmetry, which requires the action-stack $\quotstack{Y}{\cg{G}}$
        formalism and different invariants;
  \item dimensions $d\ge3$, where no analogous completeness theorem is available and where
        beyond-cohomology classes (\Cref{sec:beyond}) appear;
  \item intrinsic topological order, which is not an SPT and not covered by these indices.
\end{itemize}
Within its scope the theorem is sharp; outside it, realizability is open or obstructed, and
conflating the two is precisely the over-extrapolation the program's editorial stance
forbids.
\end{remark}

\section{The obstruction: the Kapustin--Fidkowski no-go}
\label{sec:nogo}

We now state the hard wall. It is the single result that most sharply bounds realizability,
and it is about the commuting-projector class specifically.

\begin{theorem}[Commuting-projector no-go; V-C]
\label{thm:VC}
Let $H=\sum_x P_x$ be a local commuting-projector Hamiltonian on a two-dimensional lattice
--- each $P_x$ a projector of bounded range, $[P_x,P_y]=0$ --- with a gapped ground state.
\begin{enumerate}
  \item[(i)] \emph{Electric (Kapustin--Fidkowski \cite{kapustin-fidkowski}).} If $H$ is
        $U(1)$-symmetric, the zero-temperature electric Hall conductance of the ground state
        vanishes, $\sigma_H=0$.
  \item[(ii)] \emph{Thermal (Kapustin--Spodyneiko \cite{kapustin-spodyneiko-thermal}).} The
        relative chiral central charge of the ground state vanishes, $c_-=0$, with \emph{no}
        symmetry hypothesis.
\end{enumerate}
Consequently: a $U(1)$-symmetric invertible phase with $\sigma_H\ne0$ (the integer quantum
Hall states and Chern insulators) is not $\CP$-realizable, by (i); and \emph{any} invertible
phase with $c_-\ne0$ (the $E_8$ state, and every chiral phase) is not $\CP$-realizable, by
(ii):
\[
  c\in\im\real,\ \bigl(\text{$\sigma_H(c)\ne0$ with $U(1)$}\bigr)\ \text{or}\ \bigl(c_-(c)\ne0\bigr)
  \quad\Longrightarrow\quad c\notin\im\real^{\CP}.
\]
The two obstructions are logically independent: $\sigma_H$ and $c_-$ are distinct invariants
--- the bosonic $E_8$ state has $\sigma_H=0$ (no $U(1)$ charge) yet $c_-=8$ --- and part (ii)
requires no symmetry, whereas part (i) needs $U(1)$.
\end{theorem}

\begin{proof}
Both parts are cited, not reproved. Part (i) is the theorem of Kapustin--Fidkowski
\cite{kapustin-fidkowski}: a $U(1)$-symmetric local commuting-projector Hamiltonian has
vanishing zero-temperature electric Hall conductance, taken here as the locally-computable,
integer-quantized index of \cite{kapustin-sopenko-hall}. Part (ii) is the theorem of
Kapustin--Spodyneiko \cite{kapustin-spodyneiko-thermal}, who derive a Kubo-type formula for
the thermal Hall conductance of a two-dimensional lattice system and prove that the relative
chiral central charge of any local commuting-projector Hamiltonian vanishes. The stated
obstructions are the contrapositives.
\end{proof}

\begin{remark}[How absolute is the wall?]
\label{rem:wallscope}
The vanishing in \Cref{thm:VC} is for the \emph{ordinary} on-site $U(1)$ electric response
(part (i)) and for the standard relative chiral central charge (part (ii)). Recent work of
Hsin--Kobayashi \cite{hsin-kobayashi-generalized} realizes nonzero \emph{generalized} Hall
conductivities in local commuting-projector models by using symmetries that are not
expressible through on-site charge operators; this does not contradict \Cref{thm:VC} but
shows the no-go is sharp about \emph{which} response it annihilates. The invertible-phase
obstruction we use ($\sigma_H$ for $U(1)$ phases (i), $c_-$ for chiral phases (ii)) is
unaffected.
\end{remark}

\begin{remark}[Exactly what is and is not obstructed]
\label{rem:nogo}
The theorem obstructs $\CP$-realizability, not gapped-realizability. Chern insulators exist;
the integer quantum Hall effect is a laboratory fact; the $E_8$ state is a genuine gapped
phase. What \Cref{thm:VC} says is that none of these admits an \emph{exactly-solvable
commuting-projector} description: their entanglement cannot be brought to a strict
finite-depth fixed point compatible with the nonzero gravitational/electromagnetic response.
The obstruction is thus a statement about the \emph{method} (exact solvability) as much as
about the phase. This is why the honest realizability picture must always name which class of
Hamiltonians is meant.
\end{remark}

\begin{corollary}[The $\CP$-image is a proper subgroup when chiral phases exist]
\label{cor:proper}
In any $(d,G)$ for which $\mathcal{C}_{d,G}$ contains a class with nonzero Hall conductance
or nonzero chiral central charge (in particular $d=2$ with no antiunitary protection),
$\im\real^{\CP}\subsetneq\im\real$: the inclusion of \Cref{def:realizable} is strict.
\end{corollary}

\begin{proof}
By \Cref{thm:VC} such a chiral class is not in $\im\real^{\CP}$; it is in $\im\real$ because
the phase exists as a gapped lattice system (Chern insulator / $E_8$ construction). Hence the
inclusion is strict.
\end{proof}

\begin{remark}[Relation to the completeness theorem]
There is no tension with \Cref{thm:VB}: the completeness theorem is about SPT phases with
on-site finite symmetry, whose invariants are torsion group-cohomology classes with zero Hall
conductance; the obstruction is about chiral (free-part) invertible phases, which are not SPT
in that sense and lie outside the scope of \Cref{thm:VB}. The two results partition the
$d=2$ landscape cleanly: torsion SPT (every in-cohomology class $\CP$-realized by
\Cref{thm:VA}, with completeness itself open, \Cref{rem:VB2d}) versus chiral free part
(obstructed for $\CP$, open in general).
\end{remark}

\section{A status table for realizability}
\label{sec:table}

We collect the state of knowledge into \Cref{tab:status}, organized by spatial dimension and
symmetry class, separating the torsion part of the classification from the free part, and
recording for each the best available construction, the applicable obstruction, and whether
realizability is settled, constructive, obstructed-for-$\CP$, or open. The entries are the
theorems and constructions cited above; nothing in the table is claimed beyond them.

\begin{table}[htbp]
\centering
\footnotesize
\setlength{\tabcolsep}{4pt}
\renewcommand{\arraystretch}{1.4}
\begin{tabular}{@{}L{0.13\textwidth}L{0.18\textwidth}L{0.17\textwidth}L{0.21\textwidth}L{0.22\textwidth}@{}}
\toprule
\textbf{Dim.\ / class} & \textbf{Invariant (torsion / free)} & \textbf{Construction} &
\textbf{Obstruction} & \textbf{Realizability status} \\
\midrule
$d=1$, on-site finite $G$ &
$H^2(G,U(1))$ (torsion); no free part &
CGLW / cluster (\Cref{thm:VA}) &
none (no chiral in $1+1$d) &
\textbf{Settled}: $\CP$, complete (\Cref{thm:VB}) \\

$d=1$, fermionic $\Ztwo$ &
$\Ztwo$ (torsion) &
Kitaev chain (\Cref{prop:kitaev}) &
none &
\textbf{Settled}: $\CP$ at fixed point \\

$d=2$, on-site finite $G$ &
$H^3(G,\Circ)$ (torsion); no bosonic free SPT &
CGLW (\Cref{thm:VA}) &
none for torsion &
\textbf{$\CP$-realized} (\Cref{thm:VA}); completeness open (\Cref{rem:VB2d}) \\

$d=2$, fermionic SPT ($G_f{=}G{\times}\Ztwo^f$) &
supercohomology $+$ beyond (torsion; e.g.\ $\Zint_8$ for $G{=}\Ztwo$) &
Gu--Wen \cite{gu-wen-supercohomology}; Tarantino--Fidkowski \cite{tarantino-fidkowski}; Wang--Gu \cite{wang-gu-fermionic} &
none (torsion) &
\textbf{Constructive}: $\CP$ for all $2{+}1$d fermionic SPTs \\

$d=2$, no symmetry (bosonic) &
free part $\Zint$, generator $E_8$ ($c_-=8$) &
$E_8$ / $K$-matrix (gapped, non-$\CP$) &
$c_-\ne0$: Kapustin--Spodyneiko (\Cref{thm:VC}(ii)) &
\textbf{Gapped yes, $\CP$ no} \\

$d=2$, fermionic (class A) &
free part $\Zint$, Chern number &
Chern insulator (gapped, non-$\CP$) &
$\sigma_H\ne0$, $U(1)$: Kapustin--Fidkowski (\Cref{thm:VC}(i)) &
\textbf{Gapped yes, $\CP$ no} \\

$d=3$, on-site finite $G$ &
$H^4(G,U(1))$ (torsion) + beyond-cohomology &
CGLW (in-coh.); Walker--Wang (beyond) &
partial; higher chiral data &
\textbf{Constructive} (in-coh.\ + WW); open in general \\

$d=3$, beyond cohomology &
cobordism classes (Kapustin) &
Walker--Wang (\Cref{prop:ww}) &
chiral boundary non-$\CP$ (\Cref{rem:wwchiral}) &
\textbf{Constructive} (bulk $\CP$); open in general \\

general $d$, invertible SRE &
$\pi_0\IPcond_{d,G}$ (torsion $\oplus$ free) &
per-class (above) &
$\sigma_H$/$c_-$ no-go on free/chiral part &
\textbf{Conjectural} (\Cref{sec:image}) \\
\bottomrule
\end{tabular}
\caption{Realizability status by spatial dimension and symmetry class. Here ``$\CP$''
abbreviates commuting-projector realizability; ``Gapped yes, $\CP$ no'' marks classes
realizable by some gapped local Hamiltonian but obstructed for commuting projectors by
\Cref{thm:VC} --- Kapustin--Fidkowski for the electric $U(1)$ response, Kapustin--Spodyneiko
for the chiral central charge.}
\label{tab:status}
\end{table}

Two patterns organize the table. First, the \emph{torsion} part of the classification is,
wherever a construction is known, realized by commuting-projector models, and in $d=1,2$ with
on-site finite symmetry this is complete. Second, the \emph{free} (chiral) part is
gapped-realizable but never $\CP$-realizable, uniformly, by \Cref{thm:VC}. The next section
turns these two patterns into conjectures about the whole image.

\section{The realizability image: conjectures}
\label{sec:image}

The theorems above describe $\im\real$ and $\im\real^{\CP}$ in the regions where
constructions and completeness results exist. The global shape of the two images is not a
theorem today; it is the conjectural content of this paper. We state three numbered
conjectures, in decreasing order of confidence, and phrase them in the condensed language so
that they interlock with the rest of the program.

\subsection{The image is the short-range-entangled subspectrum}

\begin{conjecture}[Realizability image; V-1]
\label{conj:V1}
The image of the comparison map
\[
  \real\colon\Phases_{d,G}\to\mathcal{C}_{d,G}
\]
restricted to uniformly gapped lattice families is exactly the short-range-entangled subspectrum:
$\im\real=\SRE_{d,G}$ as a subgroup of $\pi_0\IPcond_{d,G}$. Equivalently, an invertible
abstract class is gapped-realizable if and only if it is short-range-entangled, and the
non-realized classes form an identifiable obstruction subgroup
\[
  \mathrm{Obs}_{d,G} \;=\; \pi_0\IPcond_{d,G}\big/\im\real ,
\]
which contains, in the commuting-projector image, the chiral classes of \Cref{thm:VC}.
\end{conjecture}

\begin{remark}
Conjecture~V-1 is the essential-surjectivity statement of \Cref{prop:esssurj} made precise:
it asserts that the only obstruction to realizing an invertible abstract class is that the
class actually be invertible in the physical (short-range-entangled) sense, not merely
formally invertible in the spectrum. It is supported by every case in the table (each
realized class is SRE and each SRE class with a known construction is realized), but a
proof would require the yet-unproven identification of $\mathcal{C}_{d,G}$ with
$\pi_0\IPcond_{d,G}$ from Part~IV (Conjecture~IV-2) together with a general construction
theorem, neither of which is available.
\end{remark}

\subsection{Chiral classes: gapped yes, commuting-projector no}

\begin{conjecture}[Chiral realizability; V-2]
\label{conj:V2}
Every chiral invertible class --- one with nonzero chiral central charge or nonzero Hall
conductance --- is gapped-realizable by a non-commuting uniformly gapped local Hamiltonian,
but is never $\CP$-realizable. In symbols, writing $\mathcal{C}^{\chi}_{d,G}$ for the chiral
part,
\[
  \mathcal{C}^{\chi}_{d,G}\subseteq\im\real, \qquad
  \mathcal{C}^{\chi}_{d,G}\cap\im\real^{\CP}=\{0\}.
\]
Thus the no-go of \Cref{thm:VC} is specific to the commuting-projector class and not to
gapped lattice systems: $\im\real^{\CP}$ is exactly the non-chiral (zero-Hall) part of
$\im\real$.
\end{conjecture}

\begin{remark}
The second half of Conjecture~V-2 (never $\CP$) is the theorem \Cref{thm:VC} for the
electric/thermal Hall response; the conjectural part is that this is the \emph{only}
obstruction to $\CP$-realizability, i.e.\ that every non-chiral SRE class \emph{is}
$\CP$-realizable, and that every chiral class \emph{is} gapped-realizable by some (necessarily
non-commuting) model. The forward direction is known case by case (Chern insulators, $E_8$,
coupled-wire constructions) but not as a theorem for all chiral classes at once. This
conjecture is the precise sense in which ``exactly solvable'' and ``physically realizable''
part ways.
\end{remark}

\subsection{Realizability over profinite disorder hulls}

The third conjecture is where the condensed formalism does work that ordinary topology does
not: it concerns realizability not over a point but over a profinite disorder hull, as a
descent statement.

\begin{conjecture}[Profinite-family realizability; V-3]
\label{conj:V3}
Let $\dis=Q^{\Zint^d}$ be a profinite disorder hull ($Q$ a finite local-configuration set), so
a disordered family of models is an $\dis$-point $H\in\Ham_{d,G}(\dis)$. Suppose a class
$c\in\mathcal{C}_{d,G}$ is realized over a point (a translation-invariant model). Then $c$ is
realized by a \emph{uniformly gapped} $\dis$-family --- realizability is disorder-robust ---
if and only if a descent obstruction
\[
  o(c)\in H^1_{\mathrm{cond}}\bigl(\dis,\, \underline{\mathrm{Aut}}(c)\bigr)
\]
in the relevant condensed cohomology of $\dis$ with coefficients in the (condensed)
automorphisms of the realizing model vanishes. Finite quotients of $\dis$ detect $o(c)$: the
family is uniformly-gapped-realizable iff it is realizable at every finite resolution with a
common gap bound $\gapbound$.
\end{conjecture}

\begin{remark}
Conjecture~V-3 is the realizability analogue of the descent threads that run through the
whole program (Part~I's Lieb--Robinson descent, Part~II's crossed-product functoriality,
Part~III's uniform-gap descent). Its content is that disorder-robust realizability is not a
new hard analytic problem but a cohomological one \emph{once} the point-realizability and the
uniform gap (Part~III) are in hand: the profinite probe $\dis$ matches the physical
configuration space exactly, and the sheaf condition over its finite quotients is what
``disorder-robust'' should mean. The vanishing of $o(c)$ is the obstruction to gluing
finite-resolution realizations into a genuine $\dis$-family; when it vanishes, the crossed
product $C(\dis)\rtimes\Zint^d$ of Part~II \cite{bellissard-ncg-qhe,kellendonk-tilings}
carries the disorder-averaged invariant of $c$.
\end{remark}

\subsection{The comparison square}

The three conjectures fit into one diagram, which is the realizability face of the program.
Writing $\real^{\CP}$ for the restriction to commuting-projector models and $\gamma$ for
group completion, the square
\[
\begin{tikzcd}[column sep=large, row sep=large]
\{\text{$\CP$ lattice models}\}/\Weq \arrow[r, "\real^{\CP}"] \arrow[d, hook]
  & \im\real^{\CP} \arrow[d, hook] \\
\{\text{gapped lattice models}\}/\Weq \arrow[r, "\real"']
  & \mathcal{C}_{d,G}
\end{tikzcd}
\]
commutes, the left vertical is the honest inclusion of exactly-solvable models into all
gapped models, and the right vertical is the strict inclusion of \Cref{cor:proper}.
Conjecture~V-1 identifies the bottom image with $\SRE_{d,G}$; Conjecture~V-2 identifies the
top image with the non-chiral part; Conjecture~V-3 promotes the whole square from a point to a
profinite base $\dis$.

\section{Disorder-robust realizability over profinite hulls}
\label{sec:disorder}

We expand the profinite thread, since it is the part of the realizability story that the
condensed formalism is built for. Ordinary topology treats a disordered family as a section
of a bundle over a compact metric space; condensed mathematics treats it as an $\dis$-point,
where $\dis=Q^{\Zint^d}$ is genuinely profinite and its finite quotients are finite-resolution
disorder data. The realizability question then acquires a resolution-by-resolution structure.

\begin{definition}[Finite-resolution realizability]
\label{def:finres}
Let $\dis=\varprojlim_i\dis_i$ be the inverse limit of its finite quotients. A class $c$ is
\emph{realizable at resolution $i$} if there is a uniformly gapped family
$H^{(i)}\in\Gap_{d,G}(\dis_i)$ whose fibrewise phase is $c$. It is \emph{profinitely
realizable} if there is a compatible system $(H^{(i)})_i$ with a common gap bound
$\gapbound>0$, equivalently an $\dis$-point $H\in\Gap_{d,G}(\dis)$ with fibrewise phase $c$.
\end{definition}

The content of Conjecture~V-3 is that the passage from resolutionwise realizability to
profinite realizability is governed by a single condensed-cohomological obstruction. Two
structural facts, both consequences of the companion papers, make this plausible and are
theorems modulo their stated hypotheses.

\begin{proposition}[Uniform-gap descent enables gluing, assuming Conjecture~III-2]
\label{prop:descent}
Assume Conjecture~III-2 of Part~III (uniform-gap descent on the stated stability stratum).
Suppose each $H^{(i)}$ is uniformly gapped with the \emph{same} bound $\gapbound$, and that
the stability hypotheses of Part~III (LTQO / frustration-free with uniform constants) hold
across the tower. Then the compatible system $(H^{(i)})_i$ defines an $\dis$-point of
$\Gap_{d,G}$, and the fibrewise phase is locally constant over $\dis$.
\end{proposition}

\begin{proof}
By Part~III (the uniform-gap descent, its Conjecture~III-2 on the stated stratum) a family
over $\dis$ is uniformly gapped iff each finite-resolution approximation is gapped with a
common $\gapbound$; the hypothesis is exactly this. The sheaf/limit description of Part~I
(profinite family $=$ compatible finite data) then assembles the $H^{(i)}$ into an
$\dis$-point of $\Ham_{d,G}$, which lands in $\Gap_{d,G}$ by the uniform bound. Local
constancy of the phase is the Lieb--Robinson-continuity of the invariant (Part~II) over the
connected components of $\dis$'s finite quotients.
\end{proof}

\begin{remark}[Why an obstruction can remain]
\Cref{prop:descent} assumes a \emph{common} gap bound and a \emph{compatible} system. Neither
is automatic: the finite-resolution realizations may exist with gaps $\gapbound_i\to0$ (not
uniformly gapped, excluded from $\Gap$), or may fail to be compatible across the tower
(the phase at resolution $i$ may not lift to resolution $i+1$ without a discontinuity). The
first failure is an analytic one addressed by Part~III's hypotheses; the second is the
cohomological obstruction $o(c)$ of Conjecture~V-3, living in $H^1_{\mathrm{cond}}$ of the
hull with automorphism coefficients. This is the precise place where ``realizable at every
finite resolution'' can fail to give ``realizable over the hull,'' and it is a genuinely
condensed phenomenon: it has no counterpart for a single translation-invariant model.
\end{remark}

\begin{example}[Disorder cannot create chiral realizability]
\label{ex:disorderchiral}
Disorder does not evade \Cref{thm:VC}. If $c$ is chiral ($\sigma_H(c)\ne0$), then no
$\dis$-family of commuting-projector models realizes $c$ at any resolution, because each
fibre would be a $\CP$ model with $\sigma_H=0$; averaging over $\dis$ cannot produce a
nonzero Hall conductance from fibrewise-zero data. Thus $\im\real^{\CP}$ over $\dis$ is still
non-chiral, consistent with Conjecture~V-2. Disorder can, however, change \emph{which}
non-chiral class is realized (mobility gaps, disorder-induced transitions), and that is the
content the condensed crossed-product invariant of Part~II \cite{prodan-sb} is designed to track.
\end{example}

\section{Computational verification}
\label{sec:haskell}

The accompanying Haskell package \texttt{src/bordism-realizability/} makes the two elementary
propositions executable and self-checking. It is not a proof assistant; it is a battery of
exact finite computations that would fail loudly if the propositions were misstated, together
with QuickCheck properties that sample the parameter space.

\subsection{Kitaev-chain invariant sweep}

The module \texttt{Kitaev.hs} implements the Bogoliubov--de Gennes $d$-vector
$(d_y(k),d_z(k))$ of \Cref{prop:kitaev} and computes two invariants directly from it: the
winding number $\winding$, by summing signed angle increments of $(d_y,d_z)$ around
$k\in[0,2\pi)$ on a fine grid and rounding; and the Majorana number
$\Maj=\sgn(d_z(0))\sgn(d_z(\pi))$. \texttt{Main.hs} prints an invariant table as $\mu$ sweeps
across the transition at fixed $t=\Delta_{\mathrm{p}}=1$:
\begin{quote}\small
for $|\mu|<2$ the winding is $1$ and $\Maj=-1$ (topological); for $|\mu|>2$ the winding is
$0$ and $\Maj=+1$ (trivial); the two invariants agree (mod $2$) on the gapped locus and jump
together at $|\mu|=2$.
\end{quote}
A second routine sweeps the full $(\mu,t)$ plane and reports the phase boundary as the locus
where the winding changes, recovering $|\mu|=2|t|$. The transition is \emph{detected}, not
hard-coded: the code locates the grid cell where the invariant changes and reports the
midpoint.

\subsection{Cluster-state string order}

The module \texttt{Stabilizer.hs} implements an exact stabilizer-formalism engine over the
binary symplectic (check-matrix) representation of Pauli operators with $i$-power phase
tracking. It builds the cluster-state stabilizer group $\langle K_j\rangle$ of
\Cref{def:cluster}, computes $\langle P\rangle\in\{+1,-1,0\}$ for any Pauli string $P$ by
testing membership in the group via Gaussian elimination over $\mathbb{F}_2$, and evaluates
the string order parameter $S_{a,b}$ of \Cref{prop:cluster}. It confirms $\langle
S_{a,b}\rangle=+1$, $\langle Z_j\rangle=\langle X_j\rangle=0$, and computes the
symmetry-breaking degradation $\langle U(\theta)^\ast S_{a,b}U(\theta)\rangle$ by expanding
the rotated string into its $2^m$ Pauli terms and summing exact stabilizer expectations ---
recovering $(\cos\theta)^m$ without assuming the closed form.

\subsection{QuickCheck properties}

\texttt{Properties.hs} states one property per claim of the two propositions:
\begin{itemize}
  \item \emph{Invariant quantization}: for random $(\mu,t,\Delta_{\mathrm{p}})$ off the discriminant, the
        computed winding is exactly $0$ or $1$ and equals $\tfrac12(1-\Maj)$.
  \item \emph{Phase-boundary detection}: for random $t,\Delta_{\mathrm{p}}\ne0$, the winding is $1$ iff
        $|\mu|<2|t|$, and the detected boundary matches $|\mu|=2|t|$ within grid resolution.
  \item \emph{String-order quantization}: for random even-length windows, $\langle
        S_{a,b}\rangle=+1$ in the cluster state and $\langle Z_j\rangle=0$.
  \item \emph{String-order degradation}: for random $\theta\in[0,\pi/2]$ and window length,
        the computed degraded string order equals $(\cos\theta)^m$ within tolerance and is
        monotone decreasing in $\theta$.
\end{itemize}
\texttt{Main.hs} runs the demonstrations and all properties and exits nonzero on any failure,
so the package doubles as a regression test of \Cref{prop:kitaev,prop:cluster}.

\section{Discussion}
\label{sec:discussion}

\subsection{What the condensed viewpoint adds here}

Realizability is, at bottom, a question about explicit Hamiltonians, and one might ask what
the condensed formalism contributes beyond bookkeeping. Three things. First, it gives the
right home for the \emph{disorder} thread: Conjecture~V-3 and \Cref{sec:disorder} are
statements about $\dis$-points and descent that are awkward to phrase with ordinary bundles
and natural with profinite probes. Second, it makes the comparison map $\real$ a morphism in
a category with exact derived operations, so that ``obstruction subgroup''
$\mathrm{Obs}_{d,G}=\mathcal{C}_{d,G}/\im\real$ is a well-behaved condensed object rather
than a set-theoretic quotient. Third, it aligns realizability with the rest of the tower: the
same $\Weq$, the same stabilization, the same spectrum appear in Parts~III--V, so
realizability is literally the surjectivity of a map the program already constructs, not a
separate theory.

\subsection{Limitations}

We are candid about the boundaries of the rigorous content. The three Theorems are
repackagings of cited results into the realizability language; their mathematical substance
is due to Chen--Gu--Liu--Wen, Else--Nayak, Ogata, Kapustin--Fidkowski, and
Kapustin--Spodyneiko, and we claim only the translation and the clean $\CP$-versus-gapped
framing. The two Propositions are
elementary and self-contained but concern the simplest ($1$d, torsion) corner. Everything
about the global image (Conjectures~V-1, V-2, V-3) is open and depends on Part~IV's
still-conjectural identification of $\mathcal{C}_{d,G}$ with $\pi_0\IPcond_{d,G}$. In
particular we do \emph{not} prove that every SRE class is realizable, nor that every
non-chiral class is $\CP$-realizable; we prove the obstruction (chiral $\Rightarrow$ not
$\CP$) and organize the positive evidence.

\subsection{Non-invertible order}

The entire spectrum-level story is about invertible (SRE) phases. Intrinsic topological order
(anyons, modular tensor categories, nonzero total quantum dimension) is not classified
by $\IPcond_{d,G}$ and its realizability is a different question, requiring the condensed
\emph{higher} stack of phases and defects that Part~VI flags as out of current scope.
Commuting-projector models realize a great deal of non-chiral topological order (toric code,
string-net / Levin--Wen, Walker--Wang), but the chiral topological orders (e.g.\ the
non-abelian Ising phase with $c_-=1/2$) inherit the obstruction of \Cref{thm:VC} through
their edge and are not $\CP$-realizable.

\subsection{Open directions}

The most consequential open problem is a \emph{general construction theorem}: a procedure
that, given an arbitrary SRE class in $\pi_0\IPcond_{d,G}$, outputs a uniformly gapped
lattice model realizing it (non-commuting where the class is chiral). Conjecture~V-1 asserts
such a procedure exists; the coupled-wire and network constructions are the closest existing
approximations for chiral phases. A second direction is to prove or refute the descent
obstruction picture of Conjecture~V-3 in a tractable family (e.g.\ the disordered Kitaev
chain over a Bernoulli hull), where the crossed-product invariant of Part~II is computable and
the obstruction $o(c)$ should be explicitly checkable.

\section{Conclusion}
\label{sec:conclusion}

We have separated the realizability question, which abstract bordism/homotopy classes are
produced by explicit gapped lattice models, into a constructive region, an obstructed
region, and an open region, all inside the condensed program of this series. The constructive
region is governed by two theorems: every group-cohomology class is realized by a
commuting-projector model (Theorem~V-A), and in $d=1$ with on-site finite symmetry the
operator-algebraic index is complete, so realized equals invariant (Theorem~V-B); in $d=2$ the
index exists and every in-cohomology class is realized, though completeness is open. The
obstructed region is governed by one theorem in two parts: local commuting-projector
Hamiltonians have zero electric Hall conductance under $U(1)$ (Kapustin--Fidkowski) and zero
chiral central charge (Kapustin--Spodyneiko), so every chiral class lies outside the
commuting-projector image (Theorem~V-C), and the inclusion
$\im\real^{\CP}\subsetneq\im\real$ is strict wherever chiral phases exist. The open region is the global shape of the realizability image, stated as three
numbered conjectures: that the image is exactly the short-range-entangled subspectrum
(V-1), that chiral classes are gapped-realizable but never commuting-projector realizable
(V-2), and that realizability descends over profinite disorder hulls under a
condensed-cohomological obstruction (V-3). The two elementary propositions (the Kitaev
chain winding invariant and the cluster-state string order) anchor the torsion corner
concretely and are verified by the accompanying Haskell. Throughout, the discipline that
makes the picture coherent is the distinction between realizability by \emph{some} gapped
local Hamiltonian and realizability by a \emph{commuting-projector} model; the
Kapustin--Fidkowski (electric) and Kapustin--Spodyneiko (thermal) walls are exactly the gap
between them.

\paragraph{Code availability.}
{\hbadness=10000
\begin{sloppypar}
The accompanying Haskell package (the Kitaev-chain invariant sweep and the exact
cluster-state stabilizer engine that verify \Cref{prop:kitaev,prop:cluster}) is available
at \href{https://github.com/YonedaAI/topological-phases-of-matter}{\nolinkurl{github.com/YonedaAI/topological-phases-of-matter}}
in the directory \texttt{src/bordism-realizability/}.
\end{sloppypar}}

\phantomsection
\addcontentsline{toc}{section}{References}
\bibliographystyle{plain}
\bibliography{references}

\end{document}
